% GENERATED FILE. Edit canonical lessons, then run npm run generate:books.
% canonical-source-hash: fnv1a64:21b3d5923c8a5a5a
% canonical-lessons: KA-C430-paristhiti, KA-C430-sandarbha, KA-C430-udbhavisu, KA-C430-nibandhane, KA-C430-avasyakate

\chapter{Situations and Requirements}
\label{ch:ka-situations-and-requirements}
\hlchaptermodality{430}

\begin{chapteropening}
\textbf{By the end of this chapter:} \emph{I can say a state of affairs, a circumstance, to arise, a stipulation and a requirement.}

Complete the last lesson of chapter 430: I can say a state of affairs, a circumstance, to arise, a stipulation and a requirement.
\end{chapteropening}

\section[\texorpdfstring{\kn{ಪರಿಸ್ಥಿತಿ} (paristhiti)}{ಪರಿಸ್ಥಿತಿ (paristhiti)}]{\kn{ಪರಿಸ್ಥಿತಿ} (paristhiti) — a state of affairs; a situation}

\label{lesson:KA-C430-paristhiti}



\begin{quote}
Before the new one: say the Kannada for to suggest, then the Kannada for a solution.
\end{quote}

\subsection*{You'll want to know: \kn{ಪರಿಸ್ಥಿತಿ}}
\textbf{\kn{ಪರಿಸ್ಥಿತಿ}} — \emph{paristhiti} — \textquotedblleft{}a state of affairs; a situation\textquotedblright{}.

\begin{quote}
\kn{ಪ್ರವಾಹದ} \kn{ನಂತರ} \kn{ಊರಿನ} \kn{ಪರಿಸ್ಥಿತಿ} \kn{ಕೆಟ್ಟದಾಗಿತ್ತು}.
\end{quote}

\begin{quote}
\emph{pravāhada nantara ūrina paristhiti keṭṭadāgittu.}
\end{quote}

\begin{quote}
After the flood the situation in town was bad.
\end{quote}

\kn{ಪರಿ}- (around) + \kn{ಸ್ಥಿತಿ} (state).

\begin{cousinweb}
\textbf{\kn{ಪರಿಸ್ಥಿತಿ}} (\emph{paristhiti}) is built from Sanskrit \emph{pari-}, \textquotedblleft{}around\textquotedblright{}, and \emph{sthiti}, \textquotedblleft{}standing, state\textquotedblright{} (Monier-Williams): the state of things around.
\end{cousinweb}

\begin{grammarlens}[title={What you've built}]
One more word for saying what happens if.
\end{grammarlens}

\subsection*{Your turn}
\begin{itemize}
\raggedright
  \item \emph{Say it:} \emph{paristhiti}
  \item \emph{Say it:} \emph{paristhiti}, once more
  \item \emph{Say it:} \emph{pravāhada nantara ūrina paristhiti keṭṭadāgittu.} — After the flood the situation in town was bad.
  \item \emph{Recall:} say the Kannada for to suggest, then the Kannada for a solution, then say \emph{paristhiti} again
\end{itemize}

\begin{tcolorbox}[breakable,colback=teal!4,colframe=teal!35!black,title={Before you move on}]
What does \kn{ಪರಿಸ್ಥಿತಿ} mean? (\textquotedblleft{}A state of affairs; a situation\textquotedblright{}.) Say it once more.
\end{tcolorbox}
\section[\texorpdfstring{\kn{ಸಂದರ್ಭ} (sandarbha)}{ಸಂದರ್ಭ (sandarbha)}]{\kn{ಸಂದರ್ಭ} (sandarbha) — a circumstance; an occasion}

\label{lesson:KA-C430-sandarbha}



\begin{quote}
Before the new one: say the Kannada for a solution, then the Kannada for a state of affairs.
\end{quote}

\subsection*{You'll want to know: \kn{ಸಂದರ್ಭ}}
\textbf{\kn{ಸಂದರ್ಭ}} — \emph{sandarbha} — \textquotedblleft{}a circumstance; an occasion\textquotedblright{}.

\begin{quote}
\kn{ಇಂಥ} \kn{ಸಂದರ್ಭದಲ್ಲಿ} \kn{ಏನು} \kn{ಮಾಡಬೇಕು} \kn{ಅಂತ} \kn{ನನಗೆ} \kn{ಗೊತ್ತಿಲ್ಲ}.
\end{quote}

\begin{quote}
\emph{intha sandarbhadalli ēnu māḍabēku anta nanage gottilla.}
\end{quote}

\begin{quote}
I do not know what to do in such a situation.
\end{quote}

\kn{ಆ} \kn{ಸಂದರ್ಭದಲ್ಲಿ}, on that occasion.

\begin{cousinweb}
\textbf{\kn{ಸಂದರ್ಭ}} (\emph{sandarbha}) is borrowed from Sanskrit \emph{saṃdarbha}, \textquotedblleft{}stringing together, a composition\textquotedblright{} (Monier-Williams).
\end{cousinweb}

\begin{grammarlens}[title={What you've built}]
One more word for saying what happens if.
\end{grammarlens}

\subsection*{Your turn}
\begin{itemize}
\raggedright
  \item \emph{Say it:} \emph{sandarbha}
  \item \emph{Say it:} \emph{sandarbha}, once more
  \item \emph{Say it:} \emph{intha sandarbhadalli ēnu māḍabēku anta nanage gottilla.} — I do not know what to do in such a situation.
  \item \emph{Recall:} say the Kannada for a solution, then the Kannada for a state of affairs, then say \emph{sandarbha} again
\end{itemize}

\begin{tcolorbox}[breakable,colback=teal!4,colframe=teal!35!black,title={Before you move on}]
What does \kn{ಸಂದರ್ಭ} mean? (\textquotedblleft{}A circumstance; an occasion\textquotedblright{}.) Say it once more.
\end{tcolorbox}
\section[\texorpdfstring{\kn{ಉದ್ಭವಿಸು} (udbhavisu)}{ಉದ್ಭವಿಸು (udbhavisu)}]{\kn{ಉದ್ಭವಿಸು} (udbhavisu) — to arise; to spring up}

\label{lesson:KA-C430-udbhavisu}



\begin{quote}
Before the new one: say the Kannada for a state of affairs, then the Kannada for a circumstance.
\end{quote}

\subsection*{You'll want to know: \kn{ಉದ್ಭವಿಸು}}
\textbf{\kn{ಉದ್ಭವಿಸು}} — \emph{udbhavisu} — \textquotedblleft{}to arise; to spring up\textquotedblright{}.

\begin{quote}
\kn{ಹೊಸ} \kn{ಮನೆಯಲ್ಲಿ} \kn{ನೀರಿನ} \kn{ಸಮಸ್ಯೆ} \kn{ಉದ್ಭವಿಸಿತು}.
\end{quote}

\begin{quote}
\emph{hosa maneyalli nīrina samasye udbhavisitu.}
\end{quote}

\begin{quote}
A water problem arose in the new house.
\end{quote}

Past (it …): \kn{ಉದ್ಭವಿಸಿತು}, \emph{udbhavisitu}.

Of problems and questions.

\begin{cousinweb}
\textbf{\kn{ಉದ್ಭವಿಸು}} (\emph{udbhavisu}) is built on Sanskrit \emph{udbhava}, \textquotedblleft{}existence, origin, production\textquotedblright{} (CDIAL 2039), with the Kannada verb ending -\emph{isu}.
\end{cousinweb}

\begin{grammarlens}[title={What you've built}]
One more word for saying what happens if.
\end{grammarlens}

\subsection*{Your turn}
\begin{itemize}
\raggedright
  \item \emph{Say it:} \emph{udbhavisu}
  \item \emph{Say it:} \emph{udbhavisu}, once more
  \item \emph{Say it:} \emph{hosa maneyalli nīrina samasye udbhavisitu.} — A water problem arose in the new house.
  \item \emph{Recall:} say the Kannada for a state of affairs, then the Kannada for a circumstance, then say \emph{udbhavisu} again
\end{itemize}

\begin{tcolorbox}[breakable,colback=teal!4,colframe=teal!35!black,title={Before you move on}]
What does \kn{ಉದ್ಭವಿಸು} mean? (\textquotedblleft{}To arise; to spring up\textquotedblright{}.) Say it once more.
\end{tcolorbox}
\section[\texorpdfstring{\kn{ನಿಬಂಧನೆ} (nibandhane)}{ನಿಬಂಧನೆ (nibandhane)}]{\kn{ನಿಬಂಧನೆ} (nibandhane) — a stipulation; a regulation}

\label{lesson:KA-C430-nibandhane}



\begin{quote}
Before the new one: say the Kannada for a circumstance, then the Kannada for to arise.
\end{quote}

\subsection*{You'll want to know: \kn{ನಿಬಂಧನೆ}}
\textbf{\kn{ನಿಬಂಧನೆ}} — \emph{nibandhane} — \textquotedblleft{}a stipulation; a regulation\textquotedblright{}.

\begin{quote}
\kn{ಸಾಲದ} \kn{ನಿಬಂಧನೆಗಳನ್ನು} \kn{ನಾನು} \kn{ಗಮನವಿಟ್ಟು} \kn{ಓದಿದೆ}.
\end{quote}

\begin{quote}
\emph{sālada nibandhanegaḷannu nānu gamanaviṭṭu ōdide.}
\end{quote}

\begin{quote}
I read the terms of the loan carefully.
\end{quote}

Plural: \kn{ನಿಬಂಧನೆಗಳು}, \emph{nibandhanegaḷu}.

Mostly plural: \kn{ನಿಬಂಧನೆಗಳು}.

\begin{cousinweb}
\textbf{\kn{ನಿಬಂಧನೆ}} (\emph{nibandhane}) is borrowed from Sanskrit \emph{nibandhana}, \textquotedblleft{}binding, fastening\textquotedblright{} (Monier-Williams).
\end{cousinweb}

\begin{grammarlens}[title={What you've built}]
One more word for saying what happens if.
\end{grammarlens}

\subsection*{Your turn}
\begin{itemize}
\raggedright
  \item \emph{Say it:} \emph{nibandhane}
  \item \emph{Say it:} \emph{nibandhane}, once more
  \item \emph{Say it:} \emph{sālada nibandhanegaḷannu nānu gamanaviṭṭu ōdide.} — I read the terms of the loan carefully.
  \item \emph{Recall:} say the Kannada for a circumstance, then the Kannada for to arise, then say \emph{nibandhane} again
\end{itemize}

\begin{tcolorbox}[breakable,colback=teal!4,colframe=teal!35!black,title={Before you move on}]
What does \kn{ನಿಬಂಧನೆ} mean? (\textquotedblleft{}A stipulation; a regulation\textquotedblright{}.) Say it once more.
\end{tcolorbox}
\section[\texorpdfstring{\kn{ಅವಶ್ಯಕತೆ} (avaśyakate)}{ಅವಶ್ಯಕತೆ (avaśyakate)}]{\kn{ಅವಶ್ಯಕತೆ} (avaśyakate) — a requirement; a need}

\label{lesson:KA-C430-avasyakate}



\begin{quote}
Before the new one: say the Kannada for to arise, then the Kannada for a stipulation.
\end{quote}

\subsection*{You'll want to know: \kn{ಅವಶ್ಯಕತೆ}}
\textbf{\kn{ಅವಶ್ಯಕತೆ}} — \emph{avaśyakate} — \textquotedblleft{}a requirement; a need\textquotedblright{}.

\begin{quote}
\kn{ಊರಿಗೆ} \kn{ಒಂದು} \kn{ಆಸ್ಪತ್ರೆಯ} \kn{ಅವಶ್ಯಕತೆ} \kn{ಇದೆ}.
\end{quote}

\begin{quote}
\emph{ūrige ondu āspatreya avaśyakate ide.}
\end{quote}

\begin{quote}
The town needs a hospital.
\end{quote}

What is needed takes -\kn{ದ} / -\kn{ಯ} before it.

\begin{cousinweb}
\textbf{\kn{ಅವಶ್ಯಕತೆ}} (\emph{avaśyakate}) is built on Sanskrit \emph{āvaśyaka}, \textquotedblleft{}necessary, inevitable\textquotedblright{} (Monier-Williams), with -\emph{tā}, \textquotedblleft{}-ness\textquotedblright{}.
\end{cousinweb}

\begin{grammarlens}[title={What you've built}]
One more word for saying what happens if.
\end{grammarlens}

\subsection*{Your turn}
\begin{itemize}
\raggedright
  \item \emph{Say it:} \emph{avaśyakate}
  \item \emph{Say it:} \emph{avaśyakate}, once more
  \item \emph{Say it:} \emph{ūrige ondu āspatreya avaśyakate ide.} — The town needs a hospital.
  \item \emph{Recall:} say the Kannada for to arise, then the Kannada for a stipulation, then say \emph{avaśyakate} again
\end{itemize}

\begin{tcolorbox}[breakable,colback=teal!4,colframe=teal!35!black,title={Before you move on}]
What does \kn{ಅವಶ್ಯಕತೆ} mean? (\textquotedblleft{}A requirement; a need\textquotedblright{}.) Say it once more.
\end{tcolorbox}
